\chapter{The topological side}
\label{ch:topological-k-homology}

\section{K-homology of a cocompact piece}
Let $\Gamma$ be countable and discrete. Use locally compact,
second-countable Hausdorff proper $\Gamma$-spaces with continuous actions.
The concrete universal model is the geometric telescope from Chapter 4.
Choose a common universe containing the space and module carriers.

\begin{definition}\label{def:piece-k-homology}
\uses{def:cocompact-pieces,def:equivariant-kk}
\lean{BC4lean.TopologicalKHomology.spaceGroup,BC4lean.TopologicalKHomology.pieceGroup,BC4lean.KKTheory.CZero.action,BC4lean.KKTheory.CZero.action_apply,BC4lean.TopologicalKHomology.cZero_separableSpace}
\leanok
For each closed invariant cocompact piece $S$, define its K-homology
to be the genuine full Chapter 5 group
$KK_i^\Gamma(C_0(S,\mathbb C),\mathbb C)$ for $i=0,1$.
The source action is $(g\cdot a)(x)=a(g^{-1}x)$ and the scalar
coefficient action is trivial. Separability of the actual supremum-norm
$C_0$ algebra follows by isometric extension by zero to the one-point
compactification.
\end{definition}

\section{Transition maps}
\begin{proposition}\label{prop:topological-transition}
\uses{def:piece-k-homology,prop:kasparov-product}
\lean{BC4lean.TopologicalKHomology.transition,BC4lean.TopologicalKHomology.transition_refl,BC4lean.TopologicalKHomology.transition_trans,BC4lean.KKTheory.CZero.pieceRestriction,BC4lean.KKTheory.CZero.pieceRestriction_comp,BC4lean.KKTheory.CZero.pieceRestriction_equivariant}
\leanok
For $S\subseteq T$, actual restriction $C_0(T)\to C_0(S)$ is an
equivariant nonunital complex star homomorphism. It induces the forward
additive map on K-homology by contravariance in the first KK-variable.
Identity inclusions act as identities and successive inclusions compose.
\end{proposition}
\begin{proof}
\leanok
Closed inclusions are proper and therefore preserve vanishing at infinity
under pullback. Apply the proved Chapter 5 source-pullback laws to these
restriction homomorphisms; the two successive restrictions are their actual
composite on every function.
\end{proof}

\section{The filtered colimit}
\begin{definition}\label{def:topological-side}
\uses{prop:topological-transition}
\lean{BC4lean.TopologicalKHomology.diagram,BC4lean.TopologicalKHomology.group,BC4lean.TopologicalKHomology.ofPiece_eq_iff,BC4lean.TopologicalKHomology.existsUnique_desc,BC4lean.TopologicalKHomology.cofinalIndexEquiv,BC4lean.TopologicalKHomology.pushforward_homotopy,BC4lean.TopologicalKHomology.homotopyEquiv,BC4lean.TopologicalKHomology.universalModelEquiv,BC4lean.TopologicalKHomology.universalComparison_unique,BC4lean.TopologicalKHomology.topologicalGroup,BC4lean.TopologicalKHomology.toCanonicalModelEquiv}
\leanok
Construct
\[
 K_i^{\mathrm{top}}(\Gamma)=K_i^\Gamma(\underline E\Gamma)
 :=\varinjlim_S KK_i^\Gamma(C_0(S),\mathbb C).
\]
Here $S$ ranges over all closed invariant cocompact pieces of the chosen
model. The constructed additive group satisfies the colimit universal
property; two representatives agree exactly when they agree in a larger
piece. Replacing this system by any monotone cofinal directed family
induces a canonical additive equivalence. The cofinality condition is
$\forall S\,\exists j\colon S\subseteq S_j$.

The group is independent of the locally compact second-countable genuine
universal proper CW model satisfying the Chapter 4 fixed-point criterion.
Every equivariant comparison map induces the same additive equivalence.
The canonical implementation uses the actual geometric telescope.
\end{definition}

The colimit is the quotient of the disjoint union of the actual piece groups
by eventual equality. Its group operations and universal property are proved.
Model comparison is genuinely piecewise: an equivariant map restricted to
a cocompact source piece is proper and has a closed cocompact image.
Compact orbit representatives and finite transporters prove this fact.
The same argument bounds the entire interval of a homotopy inside one
target piece. Proper homotopies then give point-norm $C_0$ homotopies and
equal maps on the full even and odd KK groups. Chapter 4's proved CW
uniqueness supplies inverse comparisons up to equivariant homotopy.

Cofinality is necessary: directed families that merely cover every point
can have different colimits. No global properness of a universal comparison
map is assumed. This chapter constructs the topological side and does not
assert that the assembly map is an isomorphism.
